\hypertarget{ux5deux5d1ux5d5ux5d0}{%
\section{מבוא}\label{ux5deux5d1ux5d5ux5d0}}

\hypertarget{ux5e8ux5e7ux5e2-ux5d4ux5d9ux5e1ux5d8ux5d5ux5e8ux5d9-ux5deux5deux5d5ux5d3ux5dcux5d9ux5dd-ux5e8ux5e7ux5d5ux5e8ux5e1ux5d9ux5d1ux5d9ux5d9ux5dd-ux5dcux5d8ux5e8ux5e0ux5e1ux5e4ux5d5ux5e8ux5deux5e8ux5d9ux5dd}{%
\subsection{רקע היסטורי: ממודלים רקורסיביים
לטרנספורמרים}\label{ux5e8ux5e7ux5e2-ux5d4ux5d9ux5e1ux5d8ux5d5ux5e8ux5d9-ux5deux5deux5d5ux5d3ux5dcux5d9ux5dd-ux5e8ux5e7ux5d5ux5e8ux5e1ux5d9ux5d1ux5d9ux5d9ux5dd-ux5dcux5d8ux5e8ux5e0ux5e1ux5e4ux5d5ux5e8ux5deux5e8ux5d9ux5dd}}

במשך עשורים, עיבוד שפה טבעית התבסס בעיקרו על מודלים רקורסיביים כגון
רשתות נוירונים רקורסיביות ורשתות זיכרון לטווח ארוך. מודלים אלו עיבדו
רצפים צעד אחר צעד, כאשר מצב נסתר h\_t עודכן בכל שלב בהתאם לקלט הנוכחי
ולמצב הקודם. עם זאת, גישה זו סבלה ממגבלות מהותיות: תלויות לטווח ארוך לא
יכלו להיות מיוצגות ביעילות עקב בעיית הכחדת הגרדיאנט, והחישוב הסדרתי מנע
מקביליות ועיבוד יעיל.

הפתרון לבעיות אלו הגיע בשנת 2017, עם פרסום המאמר המכונן של Vaswani
ועמיתיו, שהציג את ארכיטקטורת הטרנספורמר המבוססת על מנגנון קשב עצמי
בלעדי, ללא רכיב רקורסיבי כלשהו. פריצת דרך זו שינתה את פני תחום עיבוד
השפה הטבעית מן היסוד.

\hypertarget{ux5deux5d5ux5d8ux5d9ux5d1ux5e6ux5d9ux5d4-ux5d5ux5deux5d8ux5e8ux5d5ux5ea-ux5d4ux5deux5d7ux5e7ux5e8}{%
\subsection{מוטיבציה ומטרות
המחקר}\label{ux5deux5d5ux5d8ux5d9ux5d1ux5e6ux5d9ux5d4-ux5d5ux5deux5d8ux5e8ux5d5ux5ea-ux5d4ux5deux5d7ux5e7ux5e8}}

ארכיטקטורת הטרנספורמר מציעה שני יתרונות מרכזיים על פני קודמיה: ראשית,
עיבוד מקבילי מלא של כל אסימוני הרצף בו-זמנית, ושנית, יכולת לדגמן תלויות
לטווח ארוך ביעילות גבוהה. יתרונות אלו הובילו לפיתוח מודלים מאומנים מראש
כגון BERT ו-GPT-2, אשר השיגו ביצועים מרשימים במגוון רחב של משימות עיבוד
שפה טבעית.

בהקשר העברי, ישנה חשיבות מיוחדת לפיתוח מודלים המותאמים לאתגרים הייחודיים
של השפה העברית, כגון מורפולוגיה עשירה, כתיב חסר ניקוד וסדר מילים גמיש.
מחקרים כמו AlephBERT הדגימו כי אימון מודלי טרנספורמר על קורפוס עברי גדול
מוביל לשיפור משמעותי בביצועים במשימות עיבוד שפה עברית.

ספר זה נועד לספק בסיס תיאורטי ומעשי מקיף לארכיטקטורת הטרנספורמר, תוך
התמקדות מיוחדת ביישומים לעיבוד שפה עברית. המטרה היא להנגיש את הידע
המתקדם בתחום לקהל דוברי העברית ולקדם מחקר מקומי בתחום.

\hypertarget{ux5deux5d1ux5e0ux5d4-ux5d4ux5e1ux5e4ux5e8}{%
\subsection{מבנה הספר}\label{ux5deux5d1ux5e0ux5d4-ux5d4ux5e1ux5e4ux5e8}}

ספר זה מאורגן בשישה פרקים, כל אחד בונה על קודמו ומרחיב את ההבנה של
ארכיטקטורת הטרנספורמר ויישומיה.

פרק ראשון זה מציג את ההקשר ההיסטורי ואת המוטיבציה לחקר ארכיטקטורת
הטרנספורמר. פרק שני מתאר בפירוט את מרכיבי הארכיטקטורה עצמה, לרבות שכבות
הקידוד והפענוח, הקידוד המיקומי ורשתות הזנה קדימה. פרק שלישי מתמקד
במנגנוני הקשב, ובפרט בקשב עצמי ורב-ראשי, המהווים את ליבת הארכיטקטורה.

פרק רביעי סוקר מודלים מאומנים מראש כמו BERT ו-GPT-2, תוך בחינת ההבדלים
ביניהם וטכניקות הכוונון הדק. פרק חמישי מתרכז ביישומים לעיבוד שפה עברית,
ודן באתגרים הייחודיים של השפה ובפתרונות שפותחו בספרות המחקרית. פרק שישי
ואחרון מסכם את עיקרי הספר, מצביע על מגבלות קיימות ומציע כיוונים למחקר
עתידי.
